%! AP Statistics — Unit 4, Lesson 4.4 — Homework
\documentclass[10pt]{article}
\usepackage{apstats-article}
\usepackage{apstats-boxes}

\begin{document}

\pageheader{Unit 4, Lesson 4.4}{Homework}
\namedateperiod

\vspace{0.1in}

\begin{practicebox}
\small
A hospital recorded blood type and a health outcome (whether the patient
was diagnosed with a particular condition) for a sample of 240 patients.

\vspace{6pt}
\renewcommand{\arraystretch}{1.4}
\begin{tabularx}{\linewidth}{@{} l X X X X r @{}}
 & \textbf{Type A} & \textbf{Type B} & \textbf{Type AB} & \textbf{Type O} & \textbf{Total} \\
\hline
\textbf{Condition present}  & 30 & 20 & 10 & 40 & 100 \\
\textbf{Condition absent}   & 50 & 30 & 20 & 40 & 140 \\
\hline
\textbf{Total}              & 80 & 50 & 30 & 80 & 240 \\
\end{tabularx}

\vspace{8pt}
One patient is selected at random from this sample.

\vspace{8pt}
\textbf{1.} Let $A = \{\text{patient has Type A blood}\}$ and
$B = \{\text{patient has Type O blood}\}$.

\begin{enumerate}[label=(\alph*), leftmargin=2em, itemsep=6pt]
  \item Are events $A$ and $B$ mutually exclusive? Justify in context.\\
        \writeline\\[1pt]\writeline
  \item Find $P(A \cap B)$. \blank{4cm}
  \item Find $P(A)$ and $P(B)$. \blank{6cm}
\end{enumerate}

\vspace{8pt}
\textbf{2.} Let $C = \{\text{condition present}\}$ and $D = \{\text{Type B blood}\}$.

\begin{enumerate}[label=(\alph*), leftmargin=2em, itemsep=6pt]
  \item Are events $C$ and $D$ mutually exclusive? Justify in context.\\
        \writeline\\[1pt]\writeline
  \item Find $P(C \cap D)$. Show your work. \blank{5cm}
\end{enumerate}

\vspace{8pt}
\textbf{3.} Let $E = \{\text{condition present}\}$ and $F = \{\text{condition absent}\}$.

\begin{enumerate}[label=(\alph*), leftmargin=2em, itemsep=6pt]
  \item Are $E$ and $F$ mutually exclusive? Why or why not?\\
        \writeline\\[1pt]\writeline
  \item What is $P(E \cap F)$? \blank{4cm}
\end{enumerate}

\vspace{8pt}
\textbf{4.} (AP Free-Response Style) A student claims: ``Type A blood and having the
condition present are mutually exclusive events because not everyone with Type A blood
has the condition.''
\begin{enumerate}[label=(\alph*), leftmargin=2em, itemsep=6pt]
  \item Is this claim correct? Use the table to support your answer.\\
        \writeline\\[1pt]\writeline\\[1pt]\writeline
  \item Write a correct complete-sentence explanation of whether these events are
        mutually exclusive.\\
        \writeline\\[1pt]\writeline
\end{enumerate}

\vspace{8pt}
\textbf{5.} (Extension) Suppose a new event $G = \{\text{Type AB blood}\}$ is defined.
Can you find \emph{any} blood type event that is mutually exclusive with $G$?
How many can you find? Explain.\\
\writeline\\[1pt]\writeline\\[1pt]\writeline

\vspace{8pt}
\textbf{6.} (Preview) A student says: ``If two events are mutually exclusive, knowing
one occurred tells you a lot about the other.'' Does this make intuitive sense?
How might this relate to the concept of \emph{independence} (coming in Lesson 4.6)?\\
\writeline\\[1pt]\writeline\\[1pt]\writeline
\end{practicebox}

\end{document}
